\section{Evaluation}
\label{sec:eval}

\subsection{Experimental Setup}
\label{sec:eval:setup}

We evaluate the Asymmetric Virtual Page Table (AVMP) allocator on a single NVIDIA RTX 3060 12GB GPU running CUDA 13.0, PyTorch 2.12.0+cu130, and Python 3.10.19 on WSL2 Ubuntu. We execute an experimental sweep across 270 cells: 5 allocator variants, 3 synthetic workloads, 2 model specifications (\texttt{jamba\_1\_5\_mini} and \texttt{zamba2\_2\_7b}), 3 total memory pool sizes (1 GiB, 4 GiB, 8 GiB), and 3 random seeds. The entire sweep completes in 16 minutes and 14 seconds of wall time.

We enforce strict determinism by using a logical operation counter for migration throttling rather than wall-clock time. This guarantees byte-identical reproducibility across reruns for event-deterministic fields. We include \texttt{effective\_batch\_size\_p50} in the deterministic subset by construction. We exclude \texttt{goodput} and \texttt{time\_to\_first\_oom} from strict reproducibility requirements as they depend on wall-clock execution speed.

We report 95\% confidence intervals from paired bootstrap resampling for the headline claims (Table~\ref{tab:bootstrap-ci}). For each comparison we form matched pairs over the (workload, model, pool, seed) grid and resample the per-cell delta or ratio-of-means 10{,}000 times. Variants share random seeds across cells, so the comparison is paired by construction. Pre-registered RNG seed 20260520 makes the CIs byte-stable across reruns of \texttt{scripts/bootstrap\_ci.py}.

We generate synthetic workloads to capture three distinct prompt distributions. The \texttt{uniform\_short} workload simulates KV-heavy traffic with low concurrency variance. The \texttt{mixed\_long} workload simulates long context requests that generate high SSM pressure and lower batch density. The \texttt{agentic\_burst} workload tests allocator responsiveness using a mix of short and long contexts with variable arrival loads.

\paragraph{Baseline fidelity.}
We model two production systems as allocator-level baselines rather than reproducing full inference engines. The \texttt{padded\_unified} baseline implements the unified pool padding behavior documented in vLLM issue \#37121, where SSM state is padded to attention page granularity. The \texttt{fixed\_dual\_mr05} and \texttt{fixed\_dual\_mr09} baselines model SGLang's \texttt{HybridReqToTokenPool} and \texttt{HybridLinearKVPool} with the \texttt{mamba\_full\_memory\_ratio} parameter at 0.5 and the default 0.9, respectively~\cite{zheng2024sglang}. These baselines isolate allocator behavior from kernel execution and request scheduling, which our prototype does not implement. A full reproduction against vLLM main or SGLang main would require integrating AVMP into those engines, which we leave to future work.

\subsection{Baseline Comparison}
\label{sec:eval:baseline}

We compare the dynamic AVMP allocator against static dual-pool baselines \cite{zheng2024sglang} and unified pool baselines \cite{kwon2023pagedattention}. Table~\ref{tab:baseline-comparison} presents the cross-workload aggregated results.

\begin{table*}[t]
\centering
\caption{Cross-workload $N_{\mathrm{OOM}}$ baseline comparison ($\downarrow$ lower is better). Cells report $\mathrm{mean} \pm \sigma$ where $\sigma$ is the per-cell standard deviation across 3 random seeds propagated across the 12 summed cells as $\sqrt{\sum_i \sigma_i^2}$. \textbf{Bold} = best per column; \underline{underline} = second best (ranked on means). The $\Delta\%$ column reports the change in total $N_{\mathrm{OOM}}$ versus \texttt{padded\_unified}: negative numbers mark improvement.}
\label{tab:baseline-comparison}
\input{tables/generated/table_baseline_comparison.tex}
\end{table*}

\begin{figure}[t]
\centering
\includegraphics[width=\linewidth]{figures/generated/fig_oom_comparison_final.pdf}
\caption{Cross-allocator $N_{\mathrm{OOM}}$ comparison ($\downarrow$ lower is better). Grouped bars show total OOM events per workload across 4 allocator variants; each bar is summed over 12 cells (2 model specs $\times$ 3 pool budgets $\times$ 3 seeds). \texttt{avmp\_dynamic\_b128} records the lowest counts on \texttt{mixed\_long} and \texttt{agentic\_burst}, and ties the static baselines on \texttt{uniform\_short} within 1.7 OOMs (9.0 vs 7.3, both far below \texttt{padded\_unified}'s 482.7). \texttt{padded\_unified} fails worst on every workload.}
\label{fig:oom-comparison-final}
\end{figure}

The \texttt{avmp\_dynamic\_b128} variant records the lowest cross-workload OOM count at 510, a 7.6\% reduction relative to the best static baseline, \texttt{fixed\_dual\_mr05}, which records 552 OOMs. The paired bootstrap on the per-cell delta over the 36-cell grid puts the 95\% confidence interval at $[-5.83, -1.39]$ OOMs per cell, equivalent to a 3.0--12.7\% cross-workload reduction and excluding the null. The \texttt{avmp\_static\_mr05} variant perfectly ties the \texttt{fixed\_dual\_mr05} baseline at 552 OOMs (bootstrap CI $[0,0]$), validating that our virtual handle abstraction introduces no measurable overhead relative to direct pool access.

The cross-workload reduction is concentrated on the long-context and bursty workloads. The per-cell delta on \texttt{mixed\_long} has 95\% CI $[-10.3, -1.3]$ OOMs and \texttt{agentic\_burst} has $[-9.3, -1.3]$; both exclude zero. On \texttt{uniform\_short} the per-cell delta is $+0.4$ with 95\% CI $[0.0, +1.3]$, so the comparison is statistically inconclusive on that workload, consistent with the within-1.7-OOM tie shown in Figure~\ref{fig:oom-comparison-final}. AVMP wins where the workload shifts pressure between pools; on the KV-only workload it is statistically indistinguishable from the best static partition.

Conversely, the \texttt{padded\_unified} variant performs worst, triggering 1567.7 OOM events (bootstrap CI on the per-cell delta vs \texttt{fixed\_dual\_mr05}: $[+46.2, +128]$). This confirms that our modeled unified padding baseline fails under hybrid architecture constraints due to capacity overestimation, consistent with the behavior reported in vLLM issue \#37121. Furthermore, the \texttt{fixed\_dual\_mr09} variant, matching the default 0.9 ratio used in SGLang, records 1221.3 OOMs (bootstrap CI on the per-cell delta: $[+29.2, +86.9]$), demonstrating that static default ratios perform poorly on mixed workloads. As a design trade-off, AVMP maintains a 2$\times$ VRAM footprint, peaking at 9216 MiB reserved compared to 5120 MiB for static baselines.

\begin{table*}[t]
\centering
\caption{Paired bootstrap 95\% confidence intervals for the headline claims. Each row resamples matched ``(workload, model, pool, seed)'' tuples with replacement 10{,}000 times and reports the per-tuple delta or the ratio-of-means as appropriate. The ``Significant'' column marks rows whose CI excludes the null (0 for deltas, 1 for ratios). Reproducible via \texttt{python scripts/bootstrap\_ci.py}.}
\label{tab:bootstrap-ci}
\input{tables/generated/table_bootstrap_ci.tex}
\end{table*}

\subsection{Throughput Analysis}
\label{sec:eval:throughput}

We measure system throughput using goodput, defined as completed requests per second. Our pre-registered protocol dictates that dynamic allocation is justified if the goodput exceeds 1.10$\times$ the best baseline on at least one workload. AVMP passes this threshold across all three workloads simultaneously.

Table~\ref{tab:per-workload-winner} summarizes per-workload goodput. AVMP records 434.24 req/s on \texttt{uniform\_short} versus 32.65 req/s for \texttt{fixed\_dual\_mr05}, a 13.30$\times$ ratio (paired bootstrap 95\% CI on the ratio of means: $[11.18, 16.00]$). On \texttt{mixed\_long}, 65.07 vs 27.25 req/s gives 2.39$\times$ (95\% CI $[1.70, 3.04]$). On \texttt{agentic\_burst}, 46.91 vs 25.69 req/s gives 1.83$\times$ (95\% CI $[1.42, 2.60]$). All three CIs exclude both the unit ratio and the pre-registered 1.10$\times$ threshold.

\begin{table*}[t]
\centering
\caption{Per-workload goodput comparison between AVMP and the best static baseline. Goodput columns ($\uparrow$ higher is better) are mean req/s across 12 cells; the ratio column is $g_{\mathrm{AVMP}} / g_{\mathrm{baseline}}$ where $g$ is goodput. \textbf{Bold} = winner per row. AVMP wins all three workloads, with the largest gap on \texttt{uniform\_short} ($13.30\times$).}
\label{tab:per-workload-winner}
\input{tables/generated/table_per_workload_winner.tex}
\end{table*}

We explicitly report per-workload ratios rather than a cross-workload mean because static baselines exhibit high variance across distinct prompt distributions. The \texttt{fixed\_dual\_mr05} goodput ranges narrowly from 25.69 to 32.65 req/s across the workloads. Mean aggregation obscures this variance and misleads interpretations of allocator resilience. AVMP maintains consistently higher goodput by adapting dynamically.

The primary mechanism driving this throughput increase is faster recovery from OOM events, not sustained concurrency. The median effective batch size (\texttt{effective\_batch\_size\_p50}) is strictly identical across all 5 variants per workload: 129 for \texttt{agentic\_burst}, 132 for \texttt{mixed\_long}, and 284 for \texttt{uniform\_short}. Since operational batch sizes remain workload-dominated rather than allocator-dominated, the goodput improvements directly result from static baselines spending significantly more wall time blocked in OOM-rejected recovery states.

\subsection{Sensitivity Analysis}
\label{sec:eval:sensitivity}

We conduct a two-stage sensitivity analysis to identify the dominant parameters governing the dynamic rebalancing state machine.

In Stage 1, we sweep the \texttt{migration\_batch\_size} parameter across 9 values ranging from 1 to 256. The results confirm our hypothesis that migration batch size acts as the dominant performance axis. We select \texttt{b128} as the default. Cross-workload, b128 records 510.0 OOMs versus b256's 509.3 OOMs, a 0.7 OOM gap within per-cell standard deviation (0.8 to 3.0). However, b128 migrates 313 MiB versus b256's 352 MiB per cell on average, a 12.5\% reduction in migration churn. Per-workload optima vary slightly (b4 for \texttt{uniform\_short}, b128 for \texttt{mixed\_long}, b256 for \texttt{agentic\_burst}), but all values within {b64, b128, b256} cluster within 13 OOMs cross-workload. We choose b128 as the conservative point in the near-optimal region (Figure~\ref{fig:oom-vs-batch-size}). Table~\ref{tab:stage1-batchsize} details the complete Stage 1 distributions.

\begin{figure}[t]
\centering
\includegraphics[width=\linewidth]{figures/generated/fig_oom_vs_batch_size.pdf}
\caption{$N_{\mathrm{OOM}}$ variance as a function of migration batch size $B \in \{1, \ldots, 256\}$ ($\downarrow$ lower is better). Solid lines plot AVMP per workload; dotted reference lines mark the \texttt{fixed\_dual\_mr05} static baseline for the same workload.}
\label{fig:oom-vs-batch-size}
\end{figure}

\begin{table*}[t]
\centering
\caption{Stage 1 sweep over migration batch size $B$: per-workload $N_{\mathrm{OOM}}$ for $B \in \{1, \ldots, 256\}$ ($\downarrow$ lower is better). Cells report $\mathrm{mean} \pm \sigma$ where $\sigma$ is propagated across 12 cells via $\sqrt{\sum_i \sigma_i^2}$ from per-cell std across 3 seeds. \textbf{Bold} = best per column; \underline{underline} = second best (ranked on means).}
\label{tab:stage1-batchsize}
\input{tables/generated/table_stage1_batchsize.tex}
\end{table*}

In Stage 2, we evaluate trigger threshold sensitivity at the \texttt{b128} configuration. We sweep four variants combining \texttt{threshold\_high} (0.10, 0.20) and \texttt{threshold\_low} (0.02, 0.10). We pre-registered hypotheses that lower high thresholds would benefit \texttt{mixed\_long} and lower low thresholds would benefit bursty workloads. The empirical data falsifies both hypotheses. All four threshold variants result in identical OOM counts (510.0), identical rebalance event counts, and byte-identical total migrated bytes.

\begin{figure}[t]
\centering
\includegraphics[width=\linewidth]{figures/generated/fig_threshold_sensitivity.pdf}
\caption{Stage 2 threshold sensitivity ($\downarrow$ lower is better). Bars are total $N_{\mathrm{OOM}}$ across 12 cells $\times$ 3 workloads = 36 measurements for each of 4 threshold variants plus the b128 reference. All five bars land at 510.0, confirming the stage-2 null result: threshold tuning within the sampled ranges has no measurable effect on OOM count at fixed $B = 128$.}
\label{fig:threshold-sensitivity}
\end{figure}

We systematically explored the parameter space; thresholds have a marginal effect within the sampled ranges (Figure~\ref{fig:threshold-sensitivity}). Table~\ref{tab:stage2-threshold} confirms the performance invariance across threshold bounds. This negative result narrows the design space: future work on AVMP should focus on migration batch size rather than threshold tuning.

\begin{table*}[t]
\centering
\caption{Stage 2 threshold sweep at b128 ($\downarrow$ lower \texttt{total\_oom} is better): all four threshold variants and the b128 reference yield byte-identical OOM means and identical rebalance counts. The $\pm$ values on \texttt{total\_oom} propagate per-cell std across 12 cells via $\sqrt{\sum_i \sigma_i^2}$ and are also identical, confirming the tie at the distribution level.}
\label{tab:stage2-threshold}
\input{tables/generated/table_stage2_threshold.tex}
\end{table*}

\subsection{Limitations}
\label{sec:eval:limitations}

We acknowledge four primary limitations in the current AVMP prototype design.

First, the system requires a 2$\times$ VRAM footprint trade-off (Figure~\ref{fig:peak-reserved-tradeoff}). We size the backing stores to the maximum possible capacity of both pools combined. AVMP peaks at 9 GiB on a 4 GiB pool budget to enable zero-copy migration. Integrating \texttt{cuMemMap} would address this overhead by dynamically mapping virtual addresses to physical pages, similar to the approach in vTensor~\cite{xu2024vtensor}.

\begin{figure}[t]
\centering
\includegraphics[width=\linewidth]{figures/generated/fig_peak_reserved_tradeoff.pdf}
\caption{Peak reserved VRAM trade-off for dynamic allocation; lower VRAM is better in absolute terms, but AVMP intentionally trades 2$\times$ VRAM for capacity migration headroom. Bars show mean across 12 cells with error bars at $\pm 1\sigma$ (cross-cell standard deviation), and value labels report $\mathrm{mean} \pm \sigma$.}
\label{fig:peak-reserved-tradeoff}
\end{figure}

Second, we execute the allocator as a pure Python prototype without Triton or CUDA kernels. This validates allocator semantics; integration with Triton kernels and real model runtime is future work.

Third, our evaluation uses synthetic prompt distributions. While our three presets successfully capture short, long, and bursty axes, they do not replicate real model traffic semantics. Future evaluations require ShareGPT \cite{sharegpt} and Alpaca \cite{taori2023alpaca} traces.

Finally, we test strictly on a single GPU. Extending AVMP to multi-GPU environments requires implementing tensor parallelism with cross-device virtual pool sizing synchronization.

\paragraph{Failure modes.}
AVMP does not help in several allocator regimes. When both pools simultaneously approach saturation, the rebalancing trigger fires but finds no donor pool with free fraction above \texttt{threshold\_high}, so the allocator cannot rebalance and falls back to the current partition's admission behavior; in this regime AVMP provides no additional capacity advantage over a static partition. When \texttt{migration\_batch\_size} is mistuned to extreme values, the allocator either migrates too slowly to recover from CapacityError (at $B=1$) or incurs higher migration churn without proportional OOM improvement (at $B \geq 256$, see Figure~\ref{fig:oom-vs-batch-size} and Table~\ref{tab:stage1-batchsize}). Invalid configurations where \texttt{threshold\_low} $\geq$ \texttt{threshold\_high} produce undefined transition semantics; the prototype includes a validation check but does not formally prove allocator invariants. The logical operation counter uses a 64-bit integer with no wraparound risk in any realistic deployment lifetime, but distributed multi-instance deployments would require additional synchronization beyond the current single-process design.

\subsection{Reproducibility}
\label{sec:eval:reproducibility}

The full sweep harness, generated tables and figures, paper source, and pre-registered analysis protocol are available at \url{https://github.com/codepawl/cachepawl}. The 270-cell sweep completes in 16:14 wall time on a single NVIDIA RTX 3060 12GB. Event-deterministic fields (OOM counts, rebalance events, migrated bytes) reproduce byte-identically across reruns; goodput and time-to-first-OOM depend on wall-clock execution speed and are not subject to byte-identical reproducibility.
